% Preamble
\documentclass[11pt, aspectratio=169]{beamer}
	% Packages
	\usepackage[english]{babel}
	\usepackage{lipsum}
	\usepackage{mathptmx}
	\usepackage{xcolor}
	% Themes
	\usefonttheme[onlymath]{serif}
	\usetheme[nofirafonts]{focus}
	\renewcommand{\thempfootnote}{\arabic{mpfootnote}}
	% Cover
	\title{Microeconometrics Using Stata}
	\subtitle{Nonlinear optimization methods}
	\author{
		Jhon R. Ordoñez
		\footnote{\label{1} National University of San Cristóbal de Huamanga}
		\footnote{\label{2} Faculty of Economic, Administrative and Accounting Sciences}
		\footnote{\label{3} Professional School of Economics}
	}
	\date{\today}
% Body
\begin{document}
	% Cover
	\begin{frame}
		\maketitle
	\end{frame}
	% Outline
	\begin{frame}[t]
		\frametitle{Outline}
		\tableofcontents
	\end{frame}
	% Section 1
	\section{Exercise 1}
		% Slide 1
		\begin{frame}
			\frametitle{Exercise 1}
				\begin{block}{Exercise 1}
					Consider estimation of the logit model covered in \textcolor{red!50}{\textbf{section 10.5}} and
					\textcolor{red!50}{\textbf{chapter 17}}. Then $Q(\mathbf{\beta}) = \sum_{i}\left\{y_{i} \ln{\Lambda_{i}} + \left(1 - y_{i}\right)\Lambda_{i}\right\}$, where
					$\Lambda_{i} = \Lambda\left(\mathbf{x}_{i}'\mathbf{\beta}\right) = exp\left(\mathbf{x}_{i}'\mathbf{\beta}\right)/ \left\{1 + \mathbf{x}_{i}'\mathbf{\beta}\right\}$. 
					Show that $g\left(\mathbf{\beta}\right) = \sum_{i}\left(y_{i} - \Lambda_{i}\right)\mathbf{x}_{i}$
					and $H\left(\mathbf{\beta}\right) = \sum_{i} - \Lambda_{i}\left(1 - \Lambda_{i}\right)\mathbf{x}_{i}\mathbf{x}_{i}'$. 
					\textcolor{magenta!87}{Hint: \textbf{$\partial \Lambda \left(z\right) / \partial z = \Lambda\left(z\right)\left\{1 - \Lambda\left(z\right)\right\}$}}.
					Use the data on \textcolor{blue}{\textbf{docvis}} to generate the binary variable \textcolor{blue}{\textbf{d$\_$dv}} for whether there are
					any doctor visits. Using just 2002 data, as in this chapter, use \textcolor{blue}{\textbf{logit}} to
					perform logistic regression of the binary variable \textcolor{blue}{\textbf{d$\_$dv}} on \textcolor{blue}{\textbf{private}}, \textcolor{blue}{\textbf{chronic}},
					\textcolor{blue}{\textbf{female}}, \textcolor{blue}{\textbf{income}}, and an intercept. Obtain robust estimates of the standard
					errors. You should find that the coefficient of \textcolor{blue}{\textbf{private}}, for example, equals $1.27266$
					, with a robust standard error of $0.0896928$.
				\end{block}
		\end{frame}
	% Section 2
	\section{Exercise 2}
		% Slide 1
		\begin{frame}
			\frametitle{Exercise 2}
			\begin{block}{Exercise 2}
				Adapt the code of \textcolor{red!50}{\textbf{section 16.2.3}} to fit the logit model
				of exercise $1$ using \textbf{NR} iterations coded in Mata. 
				\textcolor{magenta!87}{\textbf{Hint:} In defining an $n\times 1$ column vector with entries $\Lambda_{i}$ , you may find it helpful
				to use the fact that $J(n,1,1)$ creates an $n \times 1$ vector of $1$s}.
			\end{block}
		\end{frame}
	% Section 3
	\section{Exercise 3}
		% Slide 1
		\begin{frame}
			\frametitle{Exercise 3}
			\begin{block}{Exercise 3}
				Adapt the code of \textcolor{red!50}{\textbf{section 16.5.3}} to fit the logit model of \textbf{exercise 1} using the
				\textcolor{blue}{\textbf{ml}} command method \textcolor{blue}{\textbf{lf}}.
			\end{block}
		\end{frame}
	% Section 4
	\section{Exercise 4}
		% Slide 1
		\begin{frame}
			\frametitle{Exercise 4}
			\begin{block}{Exercise 4}
				Generate $100,000$ observations from the following logit model \textbf{DGP},
				$y_{i} = 1$ if $\beta_{1} + \beta_{2}x_{i} > 0$ and $y_{i} = 0$ otherwise;
				$(\beta_{1},\beta_{2}) = (0,1)$; $x_{i} \sim N(0,1)$
				where $u_{i}$ is logistically distributed. Using the inverse transformation method,
				you can compute a draw $u$ from the logistic distribution as $u = -\ln\left\{(1-r)/r\right\}$
				, where $r$ is a draw from the uniform distribution. Use
				data from this \textbf{DGP} to check the consistency of your estimation method in
				\textbf{exercise 3} or, more simply, of the \textcolor{blue}{\textbf{logit}} command.
			\end{block}
		\end{frame}
	% Section 4
	\section{Exercise 5}
		% Slide 1
		\begin{frame}
			\frametitle{Exercise 5}
			\begin{block}{Exercise 5}
				Consider the \textbf{NLS} example in \textcolor{red!50}{\textbf{section 16.5.5}} with an exponential conditional
				mean. Fit the model using the \textcolor{blue}{\textbf{ml}} command and the \textcolor{blue}{\textbf{lfnls}} program. Also, fit
				the model using the \textcolor{blue}{\textbf{nl}} command, given in \textcolor{red!50}{\textbf{section 13.3.6}}. Verify that these
				two methods give the same parameter estimates but, as noted in the text, the
				robust standard errors differ.
			\end{block}
		\end{frame}
	% Section 6
	\section{Exercise 6}
		% Slide 1
		\begin{frame}
			\frametitle{Exercise 6}
			\begin{block}{Exercise 6}
				Continue the preceding exercise. Fit the model using the \textcolor{blue}{\textbf{ml}} command and the
				\textcolor{blue}{\textbf{lfnls}} program with default standard errors. These implicitly assume that the
				\textbf{NLS} model error has a variance of $\sigma^{2} = 1$. Obtain an estimate of 
					\begin{equation}
						s^{2} = \left(\dfrac{1}{N - K}\right) \sum\limits_{i} \left\{y_{i} - \exp\left(\mathbf{x}_{i}'\hat{\mathbf{\beta}}\right) \right\}^{2}
					\end{equation}
				, using the \textcolor{blue}{\textbf{predictnl}} postestimation
				command to obtain $\exp\left(\mathbf{x}_{i}'\hat{\mathbf{\beta}}\right)$. Then, obtain an estimate of the \textbf{VCE} by
				multiplying the stored result \textcolor{blue}{\textbf{e(V)}} by $s^{2}$. Obtain the standard error of $\hat{\beta}_{private}$,
				and compare this with the standard error obtained when the \textbf{NLS} model is fit
				using the \textcolor{blue}{\textbf{nl}} command with a default estimate of the \textbf{VCE}.
			\end{block}
		\end{frame}
	% Section 7
	\section{Exercise 7}
		% Slide 1
		\begin{frame}
			\frametitle{Exercise 7}
			\begin{block}{Exercise 7}
				Consider a Poisson regression of \textcolor{blue}{\textbf{docvis}} on the regressors \textcolor{blue}{\textbf{private}}, \textcolor{blue}{\textbf{chronic}},
				\textcolor{blue}{\textbf{female}}, and \textcolor{blue}{\textbf{income}} and the programs given in \textcolor{red!50}{\textbf{section 16.7}}. Run the \textcolor{blue}{\textbf{ml
			    model d0 d0pois}} command, and confirm that you get the same output as
				produced by the code in \textcolor{red!50}{\textbf{section 16.7.3}}. Confirm that the nonrobust standard
				errors are the same as those obtained using \textcolor{blue}{\textbf{poisson}} with default standard
				errors. Run \textcolor{blue}{\textbf{ml model d1 d1pois}}, and confirm that you get the same output as
				produced by the code in \textcolor{red!50}{\textbf{section 16.7.3}}. Run \textcolor{blue}{\textbf{ml model d2 d2pois}}, and confirm
				that you get the same output as that given in \textcolor{red!50}{\textbf{section 16.7.3}}.
			\end{block}
		\end{frame}
	% Section 8
	\section{Exercise 8}
		% Slide 1
		\begin{frame}
			\frametitle{Exercise 8}
			\begin{block}{Exercise 8}
				Adapt the code of \textcolor{red!50}{\textbf{section 16.7.3}} to fit the logit model of exercise 1 by using
				\textcolor{blue}{\textbf{ml}} command method \textcolor{blue}{\textbf{d0}}.
			\end{block}
		\end{frame}
	% Section 9
	\section{Exercise 9}
		% Slide 1
		\begin{frame}
			\frametitle{Exercise 9}
			\begin{block}{Exercise 9}
				Adapt the code of \textcolor{red!50}{\textbf{section 16.7.2}} to fit the logit model of exercise 1 by using
				\textcolor{blue}{\textbf{ml}} command method \textcolor{blue}{\textbf{lf1}} with robust standard errors reported.
			\end{block}
		\end{frame}
	% Section 10
	\section{Exercise 10}
		% Slide 1
		\begin{frame}
			\frametitle{Exercise 10}
			\begin{block}{Exercise 10}
				Adapt the code of \textcolor{red!50}{\textbf{section 16.7.3}} to fit the logit model of exercise 1 by using
				\textcolor{blue}{\textbf{ml}} command method \textcolor{blue}{\textbf{d2}}.
			\end{block}
		\end{frame}
	% Section 11
	\section{Exercise 11}
		% Slide 1
		\begin{frame}
			\frametitle{Exercise 11}
			\begin{block}{Exercise 11}
				Consider the negative binomial example given in \textcolor{red!50}{\textbf{section 16.5.4}}. Fit this same
				model by using the \textcolor{blue}{\textbf{ml}} command method \textcolor{blue}{\textbf{d0}}. Hint: See the Weibull example
				in $[R]$ \textbf{ml}.
			\end{block}
		\end{frame}
	% References
	\begin{frame}[t,allowframebreaks]
		\frametitle{References}
		\bibliography{biblio.bib}
		% Style
		\bibliographystyle{apalike}
		% Authors
		\nocite{cameron2022-I}
		\nocite{cameron2022-II}
	\end{frame}
	% Cover
	\begin{frame}
		\maketitle
	\end{frame}
\end{document}